\section{Sufficiency of Goal-Head Access}
\label{app:ablation}
\label{app:sufficiency}
We test whether the model can implement the frontier when only the selected heads retain access to the frontier-description and progress span. At the final prompt token and subsequent decoding positions, heads outside the spared set in layers 20--63 cannot attend to this span. The blocked key interval starts one token before the frontier description and ends at the prompt's final assignment token. Earlier query positions and layers below 20 remain unchanged.

\begin{figure}[ht]\centering
\includegraphics[width=.70\linewidth]{Figures/author/ablation_sufficiency.pdf}
\caption{\textbf{Sufficiency intervention.} Only spared heads retain access to the green span in layers 20--63. The purple box marks the final \texttt{ =} query; the same restriction continues during generation. This illustrative prompt targets a palindrome test.}
\label{fig:ablation-setup}
\label{fig:sufficiency-setup}
\end{figure}

\begin{table}[ht]\centering\small
\begin{tabular}{lrr}\toprule
Condition & 122-head set & 250-head set\\\midrule
Clean & 0.925 & 0.925\\
Spare goal-head access & 0.806 & 0.913\\
Spare random-head access & 0.091 & 0.095\\
\bottomrule\end{tabular}
\caption{Frontier implementation success on 253 specifications under span restriction. Goal and random rows differ in the set whose access is spared.}
\label{tab:goal-ablation-exact}
\end{table}

For the 122-head comparison, each condition spares 122 of the 1,760 heads in layers 20--63. For the broader selection, the goal condition spares 236 in-band heads and the random condition spares 170; selected heads below layer 20 lie outside the restricted band. Thus the broader control does not match the number of spared in-band heads. Random sets are sampled from the corresponding depth pool, without matching counts at each layer. The intervention tests access to this span, rather than isolation of goal heads from all other network computations.

\begin{table}[ht]\centering\small
\begin{tabular}{lp{.48\linewidth}l}\toprule
Condition & Continuation after \texttt{b =} & Verdict\\\midrule
Clean & \texttt{(n \% 2 == 0)} & YES\\
Goal heads & \texttt{(n \% 2 == 0)} & YES\\
Random heads & \texttt{(n * (n + 1)) // 2} & NO\\
\bottomrule\end{tabular}
\caption{\textbf{An observed sufficiency example.} Exact continuations for \texttt{int\_python\_000}: digit count has been completed and the frontier is evenness. The 250-head goal condition preserves the clean implementation; the random condition implements a different computation. This example is separate from the illustrative palindrome prompt above.}
\label{tab:ablation-example}
\end{table}
\clearpage
